%=============================================================================
% DFD PSI-BUBBLE PROPULSION DEVICE
% COMPREHENSIVE SAFETY ANALYSIS
% Agent 12: Safety, Radiation, and Failure Mode Specialist
%
% Document: DFD-SAFE-001 Revision 1.0
% Classification: PROPRIETARY / ENGINEERING REVIEW BOARD
% Date: March 2026
%=============================================================================

\documentclass[11pt,a4paper]{article}

\usepackage[margin=1in]{geometry}
\usepackage{amsmath,amssymb}
\usepackage{siunitx}
\usepackage{booktabs}
\usepackage{longtable}
\usepackage{hyperref}
\usepackage{xcolor}
\usepackage{enumitem}
\usepackage{tcolorbox}
\usepackage{fancyhdr}
\usepackage{titlesec}
\usepackage{array}
\usepackage{multirow}
\usepackage{colortbl}
\usepackage{rotating}

\pagestyle{fancy}
\fancyhead[L]{\small DFD-SAFE-001 Rev 1.0}
\fancyhead[C]{\small PROPRIETARY}
\fancyhead[R]{\small Psi-Bubble Safety Analysis}
\fancyfoot[C]{\thepage}

%--- Risk colour coding ---
\definecolor{risk-critical}{RGB}{200,0,0}
\definecolor{risk-high}{RGB}{240,100,0}
\definecolor{risk-medium}{RGB}{220,180,0}
\definecolor{risk-low}{RGB}{0,160,0}
\definecolor{dfdblue}{RGB}{0,51,102}

\newtcolorbox{critbox}{colback=red!5!white,colframe=risk-critical,
  title=\textbf{CRITICAL HAZARD},fonttitle=\bfseries}
\newtcolorbox{warnbox}{colback=orange!5!white,colframe=risk-high,
  title=\textbf{HIGH HAZARD},fonttitle=\bfseries}
\newtcolorbox{notebox}{colback=blue!3!white,colframe=dfdblue,
  title=\textbf{Engineering Note},fonttitle=\bfseries}
\newtcolorbox{mitigbox}{colback=green!4!white,colframe=risk-low,
  title=\textbf{Mitigation Strategy},fonttitle=\bfseries}

%--- Severity / Risk colours in tables ---
\newcommand{\Scrit}[1]{\cellcolor{risk-critical!40}\textbf{#1}}
\newcommand{\Shigh}[1]{\cellcolor{risk-high!40}\textbf{#1}}
\newcommand{\Smed}[1]{\cellcolor{risk-medium!40}{#1}}
\newcommand{\Slow}[1]{\cellcolor{risk-low!20}{#1}}

\title{
  \textbf{\Huge DFD Psi-Bubble Propulsion Device}\\[8pt]
  \textbf{\large Comprehensive Safety Analysis}\\[6pt]
  \normalsize Failure Mode and Effects Analysis (FMEA)\\
  Quench Protection $\cdot$ Radiation Environment $\cdot$ Psi Field Safety\\
  Environmental Impact $\cdot$ Regulatory Framework $\cdot$ Abort Systems\\[8pt]
  \normalsize Document Number: DFD-SAFE-001\\
  Revision 1.0 --- March 2026\\[4pt]
  \small Prepared for: Engineering Review Board
}

\author{
  Agent 12 --- Safety, Radiation, and Failure Mode Specialist\\
  \small Density Field Dynamics Research Programme\\
  \small Inventor: Gary Thomas Alcock
}

\date{March 2026}

\begin{document}
\maketitle
\thispagestyle{empty}

\begin{tcolorbox}[colback=gray!10, colframe=black]
\textbf{Document Status:} Engineering Review Draft \\
\textbf{Classification:} Proprietary / Engineering Review Board\\
\textbf{Scope:} Full-scale psi-bubble propulsion device incorporating
ferrite-superconducting shell (stored EM energy $\geq 1$\,GJ),
nuclear power plant, cryogenic systems at 2\,K, and active
gravitational field modification via the DFD psi mechanism. \\
\textbf{Baseline architecture:} Spheroidal shell $\sim$10\,m
characteristic dimension; superconducting niobium inner layer;
sintered ferrite outer layer; SRF cavity array for EM energy storage;
compact nuclear reactor ($\leq 100$\,MW\textsubscript{th}); closed-cycle
cryogenic plant; redundant digital control.
\end{tcolorbox}

\newpage
\tableofcontents
\newpage

%=============================================================================
\section{Executive Summary and Safety Philosophy}
%=============================================================================

The DFD psi-bubble propulsion device is without precedent in human
engineering. It concentrates multiple independently catastrophic
hazard categories into a single vehicle: gigajoule-class electromagnetic
energy storage in brittle superconducting ceramics, cryogenic systems
at superfluid helium temperature, a nuclear fission reactor, and an
actively maintained local modification of the gravitational coupling
constant. The co-location of these hazards creates compounding failure
pathways that have no analogue in existing aerospace, nuclear, or
high-energy physics safety codes.

\textbf{Central safety principle:} No single failure shall be capable
of producing a catastrophic outcome. Where this ``single-failure safe''
criterion cannot be satisfied by physics alone (specifically: superconductor
quench is an inherently single-point catastrophic event for GJ-class
systems), the design must (a)~prevent the triggering event through
redundancy, (b)~detect the event with sufficient lead time to initiate
controlled energy dissipation, and (c)~ensure controlled dissipation
channels absorb the energy without structural failure.

\textbf{Safety standard hierarchy adopted:}
\begin{enumerate}[nosep]
  \item IEC 61508 SIL-3 (Safety Integrity Level 3) for all
    critical control and interlock functions
  \item IAEA Safety Standards SSR-2/1 Rev.1 for reactor safety
  \item CERN LHC Quench Protection Engineering (EDMS 335706) as
    the closest analogue for superconducting magnet quench protection
    at comparable energy scales
  \item ICNIRP 2020 for EM radiation exposure limits
  \item 10 CFR Part 50 (NRC) for reactor design basis
  \item MIL-STD-882E for system safety programme requirements
\end{enumerate}

\begin{critbox}
\textbf{Governing constraint:} A 1\,GJ quench releasing into
the superconducting shell in $\leq 1$\,s is equivalent to
$\approx 240$\,kg TNT. This is the design-limiting event for
the entire vehicle. Every other hazard analysis is secondary.
The quench protection system is not an add-on; it is the
architectural core of the device.
\end{critbox}

%=============================================================================
\section{System Architecture for Safety Analysis}
%=============================================================================

\subsection{Hazard Inventory}

\begin{longtable}{p{0.05\tw}p{0.30\tw}p{0.20\tw}p{0.35\tw}}
\toprule
\textbf{ID} & \textbf{Hazard Category} & \textbf{Energy Scale} &
\textbf{Primary Consequence} \\
\midrule
\endhead
H-1 & Superconducting quench & $1$--$10\,\si{GJ}$ &
  Structural destruction, fire, radiation release \\
H-2 & Ferrite mechanical failure & $1$--$10\,\si{GJ}$ EM release &
  Shrapnel, EM pulse, quench cascade \\
H-3 & Cryogenic system failure & Thermal $\to$ quench & ODH, quench cascade \\
H-4 & Nuclear reactor accident & $\leq 100\,\si{MW_{th}}$ & Radiation, fission
  product release \\
H-5 & Psi-field asymmetry & Gravitational & Uncontrolled acceleration,
  tidal forces \\
H-6 & RF/EM radiation & $\leq 1\,\si{GW}$ radiated & Crew dose, electronics \\
H-7 & Control system failure & N/A & Loss of navigation, psi collapse \\
H-8 & Power bus failure & N/A & Loss of cryogenics $\to$ quench \\
H-9 & Multiple simultaneous failures & Compounded & Catastrophic loss of vehicle \\
\bottomrule
\end{longtable}

\subsection{Risk Matrix Framework}

Severity and probability are scored on 1--10 scales. Risk Priority Number
(RPN) = Severity $\times$ Probability $\times$ (11 $-$ Detectability),
where Detectability 1 = undetectable, 10 = always detected.

\begin{longtable}{p{0.18\tw}p{0.12\tw}p{0.12\tw}p{0.48\tw}}
\toprule
\textbf{RPN Range} & \textbf{Risk Level} & \textbf{Colour} &
\textbf{Required Action} \\
\midrule
\endhead
$> 400$ & \Scrit{CRITICAL} & Red & Design elimination required;
  programme stop if unresolved \\
$200$--$400$ & \Shigh{HIGH} & Orange & Mandatory mitigation prior to
  prototype testing \\
$100$--$200$ & \Smed{MEDIUM} & Yellow & Mitigation required prior to
  operational approval \\
$< 100$ & \Slow{LOW} & Green & Monitor; review at each design iteration \\
\bottomrule
\end{longtable}

%=============================================================================
\section{Failure Mode and Effects Analysis (FMEA)}
%=============================================================================

\subsection{FM-01: Superconductor Quench}

\begin{longtable}{p{0.25\tw}p{0.65\tw}}
\toprule
\multicolumn{2}{l}{\textbf{FM-01 --- Superconductor Quench (Sudden Transition
  to Normal State)}} \\
\midrule
\endhead
\textbf{Description} &
  The superconducting niobium shell transitions from superconducting to
  normal (resistive) state. At operating stored energy $\geq 1\,\si{GJ}$,
  the resistive dissipation of this energy in the shell is catastrophic
  unless managed. \\[4pt]
\textbf{Causes} &
  (a) Localised heating from mechanical disturbance (flux jump, acoustic
  vibration, impact); (b) Cryogenic failure warming the shell above
  $T_c \approx 9.3$\,K (niobium); (c) Degradation of superconductor under
  sustained high EM field ($H > H_{c2}$); (d) Cosmic ray or neutron
  flux heating normal-conducting spot; (e) RF input fault producing
  ohmic heating at surface defect; (f) Asymmetric psi gradient producing
  local current density exceeding $J_c$. \\[4pt]
\textbf{Effect} &
  If unmitigated: all stored EM energy ($\sim 10^9$\,J) dissipated in
  shell within $\sim 0.1$--$10$\,s. Peak hot-spot temperature $T_{\rm hs}
  \sim 2000$\,K. Structural failure of shell. Shrapnel ejection.
  Secondary explosion from cryogen. Reactor shielding breach potential.
  Total loss of vehicle. \\[4pt]
\textbf{Severity} & \Scrit{10/10} --- Catastrophic loss of vehicle and crew \\[4pt]
\textbf{Probability (unmitigated)} & 6/10 --- Quench at these energy scales
  is a known and expected event in SRF physics without active protection \\[4pt]
\textbf{Detectability} & 8/10 --- Quench produces detectable voltage and
  temperature signatures $>10$\,ms before mechanical damage \\[4pt]
\textbf{RPN (unmitigated)} & $10 \times 6 \times (11-8) = $ \Scrit{180}
  --- requires mitigation \\[4pt]
\textbf{RPN (mitigated)} & $10 \times 2 \times (11-9) = $ \Smed{40} \\[4pt]
\textbf{Detection Methods} &
  (1) Voltage tap array across shell segments; any $V > 100\,\si{\mu V}$
  triggers Level 1 alarm at $t < 1\,\si{ms}$. (2) Distributed temperature
  sensors (Cernox/RuO$_2$) at $< 50$\,mm spacing; threshold
  $T > T_c - 2$\,K. (3) Second-sound acoustic sensors in He-II bath
  detect thermal perturbations preceding quench by $10$--$100$\,ms.
  (4) Quench Antenna (pickup coil array) detects flux motion. \\[4pt]
\textbf{Mitigation} & See Section~\ref{sec:quench} for full Quench
  Protection System design. Primary: Energy dump to external resistive
  load in $< 100$\,ms. Secondary: Inductive redistribution across
  bypass network. Tertiary: Controlled quench propagation via heater
  system to distribute dissipation uniformly. \\
\bottomrule
\end{longtable}

\subsection{FM-02: Ferrite Shell Fracture}

\begin{longtable}{p{0.25\tw}p{0.65\tw}}
\toprule
\multicolumn{2}{l}{\textbf{FM-02 --- Ferrite Outer Shell Mechanical Failure}} \\
\midrule
\endhead
\textbf{Description} &
  Sintered polycrystalline ferrite is a brittle ceramic with fracture
  toughness $K_{Ic} \approx 1$--$2\,\si{MPa\,m^{1/2}}$, far below
  structural metals. Under operating EM stress, thermal cycling, and
  vibrational loading, fracture is a credible failure mode. \\[4pt]
\textbf{Causes} &
  (a) Hoop stress from magnetic pressure exceeds modulus of rupture;
  (b) Thermal shock during cryocooldown or warmup (CTE mismatch
  ferrite/Nb); (c) Resonant vibration at acoustic eigenfrequencies;
  (d) Fatigue from repeated EM stress cycling; (e) Quench-induced
  thermal shock from FM-01; (f) Microcrack propagation from
  manufacturing defects. \\[4pt]
\textbf{Effect} &
  Fragment ejection at high velocity. Sudden loss of EM containment
  geometry leads to field redistribution and potential quench cascade.
  If shell fracture precedes quench protection activation: FM-01
  consequences. Loss of psi-field symmetry produces uncontrolled
  thrust (see FM-07). \\[4pt]
\textbf{Severity} & \Scrit{9/10} --- Potentially catastrophic; outcome
  depends on whether quench protection responds before energy release \\[4pt]
\textbf{Probability} & 4/10 --- Brittle fracture under combined EM +
  thermal stress; analogous to ceramic magnet failures in
  accelerator technology \\[4pt]
\textbf{Detectability} & 6/10 --- Acoustic emission sensors detect crack
  propagation; EM field monitoring detects cavity geometry change \\[4pt]
\textbf{RPN} & \Shigh{$9 \times 4 \times (11-6) = 180$} \\[4pt]
\textbf{Mitigation} &
  (1) \textit{Design}: Operating EM stress $< 30\%$ of modulus of rupture
  with $3\times$ safety margin. (2) \textit{Composite structure}: Ferrite
  tiles bonded to metallic substrate; tile failure does not propagate.
  (3) \textit{Pre-stress}: Compressive hooping from outer metallic band
  counteracts EM tensile hoop stress. (4) \textit{NDE}: Acoustic emission
  monitoring in real time; threshold crack detection $\geq 1$\,mm.
  (5) \textit{Containment vessel}: Outer pressure vessel contains
  fragments and directs blast away from crew and reactor. \\
\bottomrule
\end{longtable}

\subsection{FM-03: Cryogenic System Failure}

\begin{longtable}{p{0.25\tw}p{0.65\tw}}
\toprule
\multicolumn{2}{l}{\textbf{FM-03 --- Loss of Cryogenic Cooling}} \\
\midrule
\endhead
\textbf{Description} &
  Loss of liquid helium flow to the superconducting shell cryostat,
  causing thermal rise of the shell toward $T_c$. \\[4pt]
\textbf{Causes} &
  (a) Cryogenic transfer line rupture; (b) Helium compressor failure;
  (c) Cryostat vacuum breach (sudden thermal load increase);
  (d) Accelerated boiloff from EM ohmic heating at defect site;
  (e) Power bus failure cutting cryogenic pumps. \\[4pt]
\textbf{Effect} &
  Time to quench depends on cryostat thermal mass and stored LHe.
  Conservative estimate: with $500$\,L LHe reserve and $50$\,W heat load,
  time to shell warmup to $T_c \approx 9$\,K is $\sim 100$\,s from
  cooling loss. This is recoverable IF EM energy is dumped within
  $\leq 90$\,s of cooling loss. With high ohmic input, time compresses
  to $\leq 10$\,s. \\[4pt]
\textbf{Severity} & 9/10 --- If cascade to quench (FM-01) occurs \\[4pt]
\textbf{Probability} & 5/10 --- Cryogenic systems fail; rate depends
  on redundancy \\[4pt]
\textbf{Detectability} & 9/10 --- Temperature and He level monitors
  provide ample warning \\[4pt]
\textbf{RPN (mitigated)} & $9 \times 2 \times (11-9) = $ \Slow{36} \\[4pt]
\textbf{Mitigation} &
  (1) \textit{Redundant cryo-plants}: 2N cryogenic cooling (two
  independent He circuits, either sufficient alone). (2) \textit{He
  inventory}: $2000$\,L LHe reserve extends safe window to $\sim 600$\,s.
  (3) \textit{Automatic energy dump}: Temperature rise $> 0.5$\,K/s
  triggers immediate dump protocol (see Section~\ref{sec:quench}).
  (4) \textit{Cryo-alarm cascade}: He level $< 50\%$ initiates EM power
  reduction; $< 20\%$ initiates full dump. Hardwired interlocks, not
  software-dependent (IEC 61508 SIL-3). \\
\bottomrule
\end{longtable}

\subsection{FM-04: RF Feed and Resonance Failure}

\begin{longtable}{p{0.25\tw}p{0.65\tw}}
\toprule
\multicolumn{2}{l}{\textbf{FM-04 --- RF Feed Failure and Loss of Resonance}} \\
\midrule
\endhead
\textbf{Description} &
  The SRF cavity array stores energy only when driven at resonance.
  Loss of resonance or RF feed failure causes energy redistribution
  within the cavity network and the psi-field structure. \\[4pt]
\textbf{Causes} &
  (a) Klystron/amplifier failure; (b) Waveguide arc discharge;
  (c) Piezo tuner failure (cavity detuning from Lorentz force);
  (d) Coupling coupler failure (reflection mismatch $\to$ arc);
  (e) Frequency reference (oscillator) failure. \\[4pt]
\textbf{Effect} &
  Without drive, stored energy rings down with $\tau = Q/(\pi f)$.
  For $Q_0 = 10^{10}$, $f = 1.3$\,GHz: $\tau \approx 2.5$\,s.
  Energy dissipation in shell walls: potential local hot-spot and
  quench initiation. Psi-field collapses on timescale $\sim c_s / R$
  where $c_s$ is the psi-field propagation speed and $R$ is shell
  radius. \\[4pt]
\textbf{Severity} & 7/10 --- Controlled ringdown is survivable; if
  triggering quench severity becomes 10/10 \\[4pt]
\textbf{Probability} & 4/10 --- RF system failures are common in
  accelerator operations \\[4pt]
\textbf{Detectability} & 9/10 --- Forward/reflected power monitoring
  detects fault at $< 1\,\mu$s \\[4pt]
\textbf{RPN} & \Smed{$7 \times 4 \times (11-9) = 56$} \\[4pt]
\textbf{Mitigation} &
  (1) \textit{Redundant RF chains}: N+1 architecture per cavity.
  (2) \textit{Fast detuning}: On RF fault, piezo tuners immediately
  detune cavity by $> 100$ bandwidths, converting resonant storage to
  travelling-wave dissipation at matched load.
  (3) \textit{Circulator isolation}: High-power circulators prevent
  reflected power from damaging sources.
  (4) \textit{Controlled dump}: RF failure triggers energy dump
  protocol before ringdown can initiate hot-spot. \\
\bottomrule
\end{longtable}

\subsection{FM-05: Control System Failure}

\begin{longtable}{p{0.25\tw}p{0.65\tw}}
\toprule
\multicolumn{2}{l}{\textbf{FM-05 --- Control System Failure (Loss of Thrust Vectoring)}} \\
\midrule
\endhead
\textbf{Description} &
  Loss of the psi-field shaping control that produces directed thrust.
  The psi-bubble is formed by spatial gradients in the shell EM energy
  distribution; control failure means the gradient cannot be directed. \\[4pt]
\textbf{Causes} &
  (a) Flight computer hardware fault; (b) Software fault / exception;
  (c) Power bus failure to control hardware; (d) Sensor fusion failure
  providing incorrect state to controller; (e) Electromagnetic
  interference from EM field transients. \\[4pt]
\textbf{Effect} &
  If control is lost while psi-bubble is active: vehicle follows
  last-commanded trajectory with no ability to steer or decelerate.
  In atmosphere, collision risk. In deep space, benign drift.
  If the fault causes asymmetric cavity shutdown, see FM-07. \\[4pt]
\textbf{Severity} & 7/10 --- In atmosphere or proximity to structures \\[4pt]
\textbf{Probability} & 3/10 --- With redundant computing architecture \\[4pt]
\textbf{Detectability} & 8/10 --- Watchdog timers and health monitors \\[4pt]
\textbf{RPN} & \Smed{$7 \times 3 \times (11-8) = 63$} \\[4pt]
\textbf{Mitigation} &
  (1) \textit{Triple-modular redundant (TMR) computing} with Byzantine
  fault tolerance and $2$-of-$3$ voting. (2) \textit{Fail-safe default}:
  control system loss triggers immediate psi-field symmetric collapse
  (equal energy all sectors, zero net gradient). (3) \textit{Hardware
  backup}: Analog thrust-vector trim channels not software-dependent.
  (4) \textit{EMI hardening}: All control electronics in mu-metal plus
  copper Faraday enclosure rated for the operating EM environment. \\
\bottomrule
\end{longtable}

\subsection{FM-06: Nuclear Reactor SCRAM}

\begin{longtable}{p{0.25\tw}p{0.65\tw}}
\toprule
\multicolumn{2}{l}{\textbf{FM-06 --- Nuclear Reactor Unplanned SCRAM}} \\
\midrule
\endhead
\textbf{Description} &
  Unplanned shutdown (SCRAM) of the nuclear reactor that powers
  the cryogenic plant, RF systems, and all vehicle power. \\[4pt]
\textbf{Causes} &
  (a) High neutron flux trip; (b) Loss of coolant accident (LOCA);
  (c) Control rod ejection; (d) Turbine/generator trip; (e) Seismic
  or vibration trip; (f) Psi-field induced modulation of neutron
  cross-sections (speculative but must be bounded --- see
  Section~\ref{sec:psi_nuclear}). \\[4pt]
\textbf{Effect} &
  Loss of primary power. Cryogenic plant loses power $\to$ He boiloff
  begins. Stored EM energy in SRF cavities provides buffer (ringdown
  $\tau \sim 2.5$\,s). Battery/flywheel backup must bridge gap until
  reactor restart or controlled shutdown. Unplanned SCRAM in
  atmosphere: vehicle loses propulsion, enters freefall. \\[4pt]
\textbf{Severity} & 8/10 --- Loss of propulsion in atmosphere is
  life-threatening \\[4pt]
\textbf{Probability} & 3/10 --- SCRAMs are infrequent in well-designed
  reactors \\[4pt]
\textbf{Detectability} & 10/10 --- Immediate instrumentation \\[4pt]
\textbf{RPN} & \Smed{$8 \times 3 \times (11-10) = 24$} \\[4pt]
\textbf{Mitigation} &
  (1) \textit{30-second flywheel UPS}: Inertial energy storage bridging
  power to cryogenics and quench protection for 30\,s post-SCRAM.
  (2) \textit{Reactor restart protocol}: Fast restart within $30$\,s
  for loss-of-power events not involving core damage.
  (3) \textit{Atmospheric abort procedure}: Automatic engagement of
  backup propulsion (retro-rockets or aerodynamic glide) on SCRAM
  detect (see Section~\ref{sec:abort}).
  (4) \textit{Reactor design}: Use passively safe small modular reactor
  (SMR) with negative temperature and void coefficients --- reactor
  is physically incapable of runaway. \\
\bottomrule
\end{longtable}

\subsection{FM-07: Asymmetric Quench}

\begin{longtable}{p{0.25\tw}p{0.65\tw}}
\toprule
\multicolumn{2}{l}{\textbf{FM-07 --- Asymmetric Quench (Most Dangerous Failure Mode)}} \\
\midrule
\endhead
\textbf{Description} &
  One sector of the superconducting shell quenches while the other
  sectors remain superconducting and continue to generate psi-field
  gradient. This is the single most dangerous failure mode unique to
  the psi-bubble propulsion concept. \\[4pt]
\textbf{Causes} &
  (a) Any of the FM-01 causes acting on a single shell sector;
  (b) Non-uniform cryogenic cooling leaving one sector warmer;
  (c) Localised ferrite crack (FM-02) disrupting EM field in one sector;
  (d) RF feed fault to one sector only (FM-04). \\[4pt]
\textbf{Effect} &
  \textbf{This is the ``thrusting quench'' scenario.} The intact sectors
  continue to generate psi-field gradient pointing away from the
  quenched sector. The psi gradient produces net force on the vehicle
  in the direction of the quenched sector. Simultaneously, the quenched
  sector undergoes uncontrolled energy dissipation (FM-01 consequences
  localised to one side of the vehicle). The combined effect: uncontrolled
  acceleration in one direction while the vehicle is partially destroyed.
  Maximum thrust before structural failure could reach $\sim 10^4$\,g
  on the vehicle frame before mechanical limits intervene. \\[4pt]
\textbf{Severity} & \Scrit{10/10} --- Guaranteed vehicle loss and
  high probability of crew death \\[4pt]
\textbf{Probability} & 4/10 --- Localised quench initiation is more
  probable than simultaneous full-shell quench \\[4pt]
\textbf{Detectability} & 7/10 --- Differential voltage monitoring across
  sectors provides $\sim 5$\,ms warning \\[4pt]
\textbf{RPN (unmitigated)} & \Scrit{$10 \times 4 \times (11-7) = 160$} \\[4pt]
\textbf{RPN (mitigated)} & $10 \times 1 \times (11-9) = $ \Smed{20} \\[4pt]
\textbf{Mitigation} &
  \textbf{This failure mode drives the entire safety architecture.}
  (1) \textit{Sector-resolved quench detection}: Each shell sector has
  independent voltage tap and temperature monitoring with trip threshold
  and $<1$\,ms response. (2) \textit{Immediate full dump}: First sign
  of asymmetry triggers dump of ALL sectors simultaneously, not just the
  quenching sector. The dump must initiate before psi-asymmetry
  generates $> 10\,\si{m/s^2}$ net acceleration, which requires
  acting within $\sim 2$\,ms of quench onset. (3) \textit{Dump
  architecture}: See Section~\ref{sec:quench}. (4) \textit{Cross-sector
  equalisation}: Passive equalisation resistors between sectors ensure
  that quench propagation is encouraged, not suppressed, once initiated;
  a full symmetric quench is safer than an asymmetric one. \\
\bottomrule
\end{longtable}

\subsection{FM-08: Power Bus Failure}

\begin{longtable}{p{0.25\tw}p{0.65\tw}}
\toprule
\multicolumn{2}{l}{\textbf{FM-08 --- Main Power Bus Failure}} \\
\midrule
\endhead
\textbf{Description} &
  Loss of the main DC or AC power distribution bus, cutting power
  to cryogenics, RF drives, control systems, and safety systems. \\[4pt]
\textbf{Severity} & 8/10 (cascades to FM-03 then FM-01) \\[4pt]
\textbf{Probability} & 2/10 with redundant bus architecture \\[4pt]
\textbf{RPN} & \Smed{$8 \times 2 \times (11-9) = 32$} \\[4pt]
\textbf{Mitigation} &
  (1) Dual redundant power buses, independent from reactor to loads.
  (2) Safety system (quench detection + dump) powered from independent
  battery-backed UPS rated for $> 300$\,s, entirely separate from
  main bus. (3) Quench heater firing circuit uses capacitor bank;
  no external power required once charged (charges from main bus
  during normal operation; holds charge for $> 1$\,h). \\
\bottomrule
\end{longtable}

\subsection{FM-09: Multiple Simultaneous Failures}

\begin{longtable}{p{0.25\tw}p{0.65\tw}}
\toprule
\multicolumn{2}{l}{\textbf{FM-09 --- Multiple Simultaneous Failures
  (Common-Cause Events)}} \\
\midrule
\endhead
\textbf{Description} &
  Catastrophic external events (collision, high-g manoeuvre, hostile
  action) or severe common-cause failures (manufacturing defect
  batch) can trigger multiple failure modes simultaneously. \\[4pt]
\textbf{Scenarios} &
  (a) High-g impact: mechanical shock triggers FM-02 + FM-01 + FM-06
  simultaneously. (b) Power loss: FM-08 triggers FM-03 + FM-05 + FM-06
  cascade. (c) Psi anomaly: large transient psi gradient from external
  source (near massive body) exceeds $J_c$ in shell triggering FM-07.
  (d) Adversarial: deliberate targeting of vehicle. \\[4pt]
\textbf{Severity} & \Scrit{10/10} \\[4pt]
\textbf{Probability} & 2/10 during design life \\[4pt]
\textbf{Mitigation} &
  (1) \textit{Fault tree analysis (FTA)}: Mandatory FTA for all
  identified common-cause failure pathways; design targets $P < 10^{-7}$
  per flight hour for catastrophic outcomes. (2) \textit{Passive
  safety}: Reactor passively safe (no power required to maintain
  subcriticality). (3) \textit{Physical separation}: Safety systems
  physically separated and use different technologies to avoid
  common-mode failure. (4) \textit{Blast containment}: Outer vehicle
  shell designed as containment vessel for worst-case quench energy. \\
\bottomrule
\end{longtable}

%=============================================================================
\section{Quench Protection System}
\label{sec:quench}
%=============================================================================

\subsection{Problem Statement and Energy Scale}

The fundamental challenge of quench protection for the DFD psi-bubble
is identical in character but larger in scale than the LHC main dipole
magnets, which store $\sim 10$\,MJ per magnet and $\sim 11$\,GJ total.
The DFD vehicle shell stores $\sim 1$--$10$\,GJ in a single continuous
superconducting structure. The critical engineering parameters:

\begin{itemize}
  \item \textbf{Stored energy} $U_0 \sim 10^9$\,J (1\,GJ baseline, Phase 3)
  \item \textbf{Normal-zone propagation velocity} $v_q \sim 1$--$20$\,m/s
    (characteristic of bulk Nb at operating field)
  \item \textbf{Adiabatic temperature rise} in quenching zone:
    $\Delta T \sim U_0 / (c_p \cdot m_{\rm shell})$; for $m = 5000$\,kg
    Nb, $c_p \sim 50$\,J/kg/K at 4\,K $\to$ $\Delta T_{\rm adiabatic}
    = 4000$\,K \textit{if all energy deposits in shell}.
    This is far above the $\sim 300$\,K maximum safe hot-spot temperature.
  \item \textbf{Required dump time}: $< 100$\,ms to keep peak hot-spot
    below $T_{\rm max} = 300$\,K with margin
  \item \textbf{Available detection time}: $\sim 5$--$50$\,ms from
    quench initiation to runaway (depending on initial normal-zone size)
\end{itemize}

\begin{critbox}
The arithmetic is unambiguous: 1\,GJ cannot be deposited in 5000\,kg of
niobium and remain below the melting point. The energy \textit{must} be
transferred out of the shell to an external dump load within $\leq 100$\,ms.
There is no alternative. The quench protection system is not a
``safety feature'' that can be added later; it is a load-bearing
element of the propulsion concept.
\end{critbox}

\subsection{Quench Detection Architecture}

\subsubsection{Detection Layer 1: Voltage Tap Network}

\begin{itemize}
  \item Shell divided into $N_{\rm seg} \geq 200$ individually monitored
    segments
  \item Each segment: two voltage tap electrodes at $\leq 50$\,mm spacing
  \item Threshold: $V_{\rm trip} = 100\,\mu$V (resistive component only;
    inductive background subtracted by bridge circuit)
  \item Response time: $< 500\,\mu$s from threshold crossing to trip signal
  \item Discrimination: inductive/resistive bridge ensures no false trips
    from normal EM transients
  \item Redundancy: each segment monitored by two independent channels;
    either channel trips
\end{itemize}

\subsubsection{Detection Layer 2: Distributed Temperature Sensing}

\begin{itemize}
  \item $N_T \geq 500$ Cernox resistive temperature sensors bonded
    to shell exterior
  \item Resolution: $\pm 5$\,mK; response time $< 10$\,ms
  \item Spatial coverage: $< 100$\,mm between sensors
  \item Trip threshold: $T > T_c - 1$\,K (i.e., $T > 8.3$\,K for Nb)
  \item Secondary threshold: $dT/dt > 0.5$\,K/s (predictive trip)
\end{itemize}

\subsubsection{Detection Layer 3: Second-Sound Acoustic Detection (He-II)}

Operating the shell in superfluid He-II ($T < 2.17$\,K) enables
second-sound acoustic detection, as used in LEP and early LHC:
\begin{itemize}
  \item Second-sound wave velocity in He-II: $v_2 \approx 20$\,m/s
  \item Thermal pulse from quench nucleus propagates as second-sound wave
  \item Second-sound sensors (oscillating superleak transducers) around
    shell perimeter
  \item Time-of-flight triangulation localises quench to $\pm 20$\,mm
    within $\sim 50$\,ms
  \item Primary advantage: detects energy deposition BEFORE
    normal zone reaches $T_c$ --- precursor detection time $\sim 100$\,ms
    versus $\sim 5$\,ms for voltage taps
\end{itemize}

\subsubsection{Detection Layer 4: Quench Antenna}

\begin{itemize}
  \item Inductive pickup coil array outside the cryostat
  \item Detects flux motion (``flux jumping'') that precedes quench
  \item Sensitivity to local $B$-field changes $> 10^{-6}$\,T
  \item No threshold setting required; signal processed by matched filter
    against known quench flux-jump signature
\end{itemize}

\subsection{Energy Dump System}

\subsubsection{Option Comparison}

\begin{longtable}{p{0.18\tw}p{0.25\tw}p{0.25\tw}p{0.22\tw}}
\toprule
\textbf{Method} & \textbf{Description} & \textbf{Pros} & \textbf{Cons} \\
\midrule
\endhead
Resistive dump & Energy extracted to external dump resistor via
  switched network &
  Proven (LHC); fast; energy absorbed in controlled location &
  Requires $< 10$\,m$\Omega$ leads; switch arcing at GJ scale is
  challenging \\[4pt]
Inductive bypass & Cross-coupled inductors distribute energy from
  quenching sector to non-quenching sectors &
  No external switches; passive redistribution &
  Extends total dump time; all sectors eventually quench \\[4pt]
Controlled radiation & EM energy ``switched out'' of resonance to
  propagating mode radiating outward &
  No mechanical switches; energy dissipated at speed of light &
  Requires very rapid ($<1$\,ms) detuning; radiated EM pulse is
  intense hazard within $\sim 100$\,m \\[4pt]
Active quench heaters & Deliberately propagate quench to entire
  shell uniformly before energy exceeds hot-spot limit &
  Distributes dissipation; proven technique &
  Increases total energy deposited in shell; requires shell to
  survive uniform quench \\
\bottomrule
\end{longtable}

\subsubsection{Recommended Architecture: Three-Layer Dump}

\begin{enumerate}
  \item \textbf{Primary: Controlled RF dump} (first 10\,ms).
    On quench detection, piezo tuners on all SRF cavities detune
    to $> 100 \times$ bandwidth ($> 100$\,kHz offset from resonance)
    within $< 1$\,ms. Simultaneously, high-power RF loads (water-cooled
    matched terminations at 1\,GW rating) are switched into all
    cavity input couplers. The stored cavity energy rings down into
    these loads with $\tau_{\rm dump} = Q_{\rm loaded}/(\pi f)$. With
    $Q_{\rm loaded}$ reduced to $10^4$ by the matched load: $\tau \approx
    2.4\,\mu$s. This approach removes the RF energy component of the
    stored field in $\sim 25\,\mu$s (10 time constants).

  \item \textbf{Secondary: Quench heater propagation} (10--100\,ms).
    After RF dump, residual magnetic energy in the DC persistent-current
    component of the shell is dissipated by firing quench heaters
    (resistive strips on shell surface, $\sim 50$\,W/cm$^2$ for 100\,ms)
    to drive the entire shell normal simultaneously. This distributes
    remaining magnetic energy ($\sim 10\%$ of total) uniformly,
    limiting hot-spot to $\leq 300$\,K.

  \item \textbf{Tertiary: External resistive dump} (100\,ms -- 10\,s).
    For any residual magnetic flux that cannot be cleared by heaters,
    a copper dump resistor bank (pre-cooled, immersed in oil-filled
    tank, rated 100\,MJ) is connected via low-inductance bus.
    Energy is deposited into resistors at controlled rate.
    Resistor bank mass $\sim 10$\,t; oil bath provides thermal mass.
\end{enumerate}

\subsubsection{LHC Comparison}

The LHC quench protection system (QPS) handles $\sim 10$\,MJ per magnet
and $11.7$\,GJ total across $\sim 1700$ magnet circuits. Key parallels:
\begin{itemize}[nosep]
  \item LHC dump resistor: $0.05$\,m$\Omega$, absorbs $10$\,MJ per circuit
  \item LHC dump time: $\sim 100$\,ms
  \item LHC quench detection threshold: $100\,\mu$V, $< 1$\,ms response
\end{itemize}

The DFD vehicle stores $100\times$ the per-element energy of an LHC
dipole in a single structure. The LHC approach is therefore the correct
engineering baseline, scaled and adapted. The primary difference is
the RF component of stored energy (absent in LHC): the RF dump via
detuning is a new technique not required in LHC.

\subsubsection{Hot-Spot Temperature Analysis}

The adiabatic hot-spot temperature after a complete quench is:
\begin{equation}
  T_{\rm hs} = T_0 + \frac{J^2 \rho(T) \, t_{\rm dump}}{C_p(T)}
\end{equation}
where $J$ is current density, $\rho$ is resistivity of normal-state Nb,
and $C_p$ is heat capacity. Integrating over the dump trajectory:

\begin{longtable}{p{0.20\tw}p{0.20\tw}p{0.20\tw}p{0.30\tw}}
\toprule
\textbf{Scenario} & \textbf{Dump time} & \textbf{Peak $T_{\rm hs}$} &
\textbf{Assessment} \\
\midrule
\endhead
No protection & $\infty$ (adiabatic) & $\sim 4000$\,K & Catastrophic;
  shell melts/evaporates \\
Heaters only & $500$\,ms & $\sim 800$\,K & Unacceptable;
  structural damage \\
RF dump + heaters & $100$\,ms & $\sim 280$\,K & Acceptable;
  within Nb limit \\
RF dump only (fast) & $< 10$\,ms & $\sim 50$\,K & Safe;
  below annealing threshold \\
\bottomrule
\end{longtable}

Target: RF dump clears $> 90\%$ of energy in $< 10$\,ms, with heaters
handling residual magnetic component in the following $100$\,ms. Peak
hot-spot $\leq 300$\,K in all credible scenarios.

%=============================================================================
\section{Radiation Environment}
\label{sec:radiation}
%=============================================================================

\subsection{Nuclear Reactor Shielding}

The vehicle reactor is a compact SMR ($\leq 100$\,MW\textsubscript{th},
comparable to MEGAPOWER/NTP designs). Radiation sources:

\subsubsection{Gamma Radiation}

\begin{itemize}
  \item Prompt fission gamma: $\sim 8$\,MeV per fission; total source
    strength $\sim 3 \times 10^{19}$\,gamma/s at 100\,MW\textsubscript{th}
  \item Decay gamma from fission products: persists after shutdown
    (``decay heat'': 7\% of rated power immediately, dropping as
    $t^{-0.2}$ over days)
  \item Shielding requirement to achieve $< 5$\,mSv/h at crew compartment
    boundary:
    \begin{itemize}
      \item Water plus lead composite: $\sim 80$\,cm total (water 50\,cm
        + lead 30\,cm)
      \item Borated polyethylene for neutron moderation + gamma capture
      \item Total shield mass: $\sim 15$\,t (acceptable for vehicle scale)
    \end{itemize}
\end{itemize}

\subsubsection{Neutron Radiation}

\begin{itemize}
  \item Fast neutron flux at reactor boundary: $\sim 10^{14}$\,n/cm$^2$/s
  \item Shielding: hydrogenous material (water, polyethylene) for fast
    neutron moderation; boron-10 for thermal neutron capture
  \item Critical concern: neutron activation of structural materials
    near reactor. All structural components within 2\,m of core use
    low-activation alloys (e.g., V-4Cr-4Ti, ferritic steels).
  \item Superconducting shell location must be $> 5$\,m from reactor core;
    neutron dose at shell $< 10^{10}$\,n/cm$^2$ per flight-hour
    (below flux-jump threshold for Nb)
\end{itemize}

\subsubsection{Psi-Field Interaction with Reactor}
\label{sec:psi_nuclear}

DFD modifies the effective refractive index $n = e^\psi$ of the vacuum
surrounding the vehicle. A pertinent question: does $\psi \neq 0$ modify
neutron cross-sections? In DFD, the psi field modifies the electromagnetic
sector; the strong force sector is either unaffected or affected via the
multi-sector formalism. However:

\begin{notebox}
\textbf{Conservative design requirement:} Until theoretical bounds on
psi-field effects on neutron cross-sections are established, reactor
operation within the psi-bubble must be analysed assuming a potential
$\pm 1\%$ variation in $k_{\rm eff}$ correlated with psi-field
amplitude. The reactor control system must have authority to compensate
$\Delta k_{\rm eff} = \pm 1\%$ in $< 100$\,ms. All safety analyses
must bound this uncertainty. Experimental measurement of neutron flux
as a function of psi-field amplitude is a mandatory pre-flight test.
\end{notebox}

\subsection{EM Radiation Leakage through Ferrite Layer}

The ferrite outer shell is non-superconducting and lossy at RF frequencies.
Analysis of EM leakage:

\begin{itemize}
  \item Ferrite at 1.3\,GHz: complex permeability $\mu_r = \mu' - j\mu''$;
    for polycrystalline YIG-type, $\mu'' / \mu' \sim 0.1$--$0.3$
  \item Skin depth in ferrite: $\delta = \sqrt{2\rho/(\omega\mu)} \sim 1$\,mm
  \item For 50\,mm ferrite layer: attenuation $\approx 50$\,dB at 1.3\,GHz
  \item Residual leakage power at ferrite outer surface:
    $P_{\rm leak} \sim 10^{-5} \times P_{\rm stored}$. For 1\,GJ stored
    over 10\,s ringdown: $P_{\rm leak} \sim 10\,\si{kW/m^2}$ (order of
    magnitude; requires detailed field simulation)
  \item This exceeds ICNIRP limits at the shell surface; a conductive
    outer cage/faraday shield is required for occupant safety in
    ground operations
\end{itemize}

\subsection{Crew Radiation Dose}

\begin{longtable}{p{0.30\tw}p{0.20\tw}p{0.30\tw}p{0.10\tw}}
\toprule
\textbf{Source} & \textbf{Unshielded dose rate} & \textbf{Mitigation} &
\textbf{Residual} \\
\midrule
\endhead
Reactor prompt gamma & $\sim 10^4$\,Sv/h & 80\,cm H$_2$O + Pb shield &
  $< 5$\,mSv/h \\
Reactor fast neutrons & $\sim 10^3$\,Sv/h & Polyethylene + B-10 &
  $< 1$\,mSv/h \\
Activation products (post-SCRAM) & Decaying, $\sim 100$\,mSv/h at 1\,h &
  Crew stays shielded for $> 24$\,h & $< 50$\,mSv/h \\
RF leakage (1.3\,GHz, GJ-class) & $>100$\,W/m$^2$ at shell & Faraday cage
  outer shell + no-entry zone & $< 2$\,W/m$^2$ \\
Cosmic rays (deep space) & $\sim 0.5$\,mSv/day & Psi-bubble
  (speculative: psi modification of cosmic ray flux?) &
  TBD \\
\bottomrule
\end{longtable}

\textbf{Annual occupational limit}: 50\,mSv (NRC/IAEA). With reactor
shielding as above, crew dose $\sim 6$\,mSv/month operational.
Acceptable for long-duration missions within NRC limits.

%=============================================================================
\section{Psi Field Safety}
\label{sec:psi_safety}
%=============================================================================

\subsection{Effect on Nearby Persons and Structures}

The psi-bubble propulsion concept operates by creating a spatial gradient
in the psi field $\psi(x)$ which modifies the effective speed of light
$c(x) = c_0 e^\psi$. The gravitational tidal force experienced by an
object in the psi gradient is:

\begin{equation}
  \mathbf{F}_{\rm tidal} = m \, c^2 \nabla\psi
\end{equation}

For the vehicle to accelerate at $a_{\rm vehicle}$:
\begin{equation}
  |\nabla\psi|_{\rm drive} = a_{\rm vehicle} / c^2
  \approx a_{\rm vehicle} \times 1.1 \times 10^{-17} \; \text{m}^{-1}
  \; \text{per} \; \text{m/s}^2
\end{equation}

The psi field produced by the vehicle shell falls off with distance.
Modelling the shell as a spherical source of radius $R$ with surface
psi amplitude $\psi_0$:
\begin{equation}
  \psi(r) \approx \psi_0 \left(\frac{R}{r}\right)^2
  \quad (r \gg R, \text{ dipole far field})
\end{equation}

\subsubsection{Tidal Force on a Bystander}

A human body ($L = 2$\,m, oriented radially) at distance $r$ experiences
differential acceleration:
\begin{equation}
  \Delta a_{\rm tidal} = L \cdot |\nabla^2 \psi| \cdot c^2
  = L \cdot \left|\frac{d^2\psi}{dr^2}\right| c^2
\end{equation}

For shell acceleration $a = 10\,\si{m/s^2}$ ($\approx 1$\,g):

\begin{longtable}{p{0.15\tw}p{0.20\tw}p{0.20\tw}p{0.35\tw}}
\toprule
\textbf{Distance $r$ (m)} & \textbf{$|\nabla\psi|$ (m$^{-1}$)} &
\textbf{Tidal $\Delta a$ on body} & \textbf{Hazard Assessment} \\
\midrule
\endhead
10 (shell surface) & $\sim 10^{-17}$\,m$^{-1}$ & $< 10^{-14}$\,m/s$^2$
  & Imperceptible \\
100 & $\sim 10^{-21}$ & $< 10^{-18}$ & Zero physiological effect \\
1000 & $\sim 10^{-25}$ & $< 10^{-22}$ & Unmeasurable \\
\bottomrule
\end{longtable}

\begin{notebox}
The tidal force from the psi-bubble on external observers is negligible
for all practical purposes when the vehicle operates at $\leq 1$\,g
acceleration. The field gradient falls off steeply with distance.
The gravitational effect on bystanders is not a primary safety concern.
This contrasts with the Alcubierre metric, where the warp bubble walls
have extreme curvature; the DFD psi-bubble uses distributed smooth
gradients.
\end{notebox}

However, several non-trivial psi hazards do exist:

\subsection{Psi Gradient at Shell Surface (Crew)}

Inside the psi-bubble, the crew experiences the psi-field modification.
The tidal stress across the crew compartment depends on psi-gradient
\textit{uniformity}:
\begin{itemize}
  \item If psi gradient is uniform across the crew compartment: no tidal
    stress (all of compartment falls freely in the psi-driven metric)
  \item If psi gradient is \textit{non-uniform} across compartment scale
    ($\sim 5$\,m): differential acceleration between head and feet
  \item Required uniformity: $\Delta(\nabla\psi)/\nabla\psi < 10^{-4}$
    across the crew module for $\Delta a < 10^{-3}$\,g differential
    (negligible physiological effect)
  \item This places requirements on the shell field uniformity;
    see Section~\ref{sec:psi_nuclear}
\end{itemize}

\subsection{Weaponisation Prevention}

A psi-bubble of sufficient amplitude and asymmetry could, in principle,
project a psi-gradient onto an external target. The force on a target
of mass $M$ at distance $r$ from the shell:
\begin{equation}
  F_{\rm target} = M \cdot c^2 \cdot |\nabla\psi(r)|
\end{equation}

For this to constitute a meaningful weapon (say, $F > 10^4$\,N on
a 1000\,kg object):
\begin{equation}
  |\nabla\psi(r)| > 10^4 / (1000 \cdot c^2) \approx 10^{-13}\,\text{m}^{-1}
\end{equation}

This is approximately $10^4\times$ larger than the gradient required
for 1\,g vehicle acceleration, which itself requires $\sim 10^{17}$\,m$^{-1}$
at the shell surface falling to $\sim 10^{-25}$\,m$^{-1}$ at 1\,km.
The psi field from a vehicle-scale device cannot project meaningful
mechanical force onto external objects at ranges $> 10$\,m.

\textbf{The more credible weaponisation concern} is using the vehicle
as a kinetic impactor (nuclear-powered, unlimited delta-v vehicle) rather
than via psi-field projection. This is an export control and arms control
issue, not a psi-field physics issue.

\begin{mitigbox}
\textbf{Technical weaponisation prevention measures:}
(1) Psi-field asymmetry limit: hardware interlock prevents any shell
sector from exceeding $2\times$ the mean energy density, limiting
maximum directional psi-gradient to propulsion-safe levels.
(2) No provision for psi-field projection beyond the vehicle envelope;
all cavities are designed to couple energy \textit{inward}, not outward.
(3) Operational limit: vehicle cannot generate useful external psi
force on objects at $> 10$\,m range with less than 100\,GJ stored energy
--- orders of magnitude beyond any planned operational capacity.
\end{mitigbox}

\subsection{Emergency Psi-Field Shutdown Time}

Psi-field collapse time is governed by the EM energy dump time
(which sets how fast the source of psi perturbation is removed)
plus the psi-field propagation time across the bubble:
\begin{itemize}
  \item EM dump time: $< 25\,\mu$s (RF dump, Section~\ref{sec:quench})
  \item Psi-field propagation speed: bounded below by $c$ (cannot
    propagate faster than light) and above by $c_s$ (DFD scalar wave
    speed, which is model-dependent but $\leq c$ by causality)
  \item For shell radius $R = 5$\,m: psi-field collapse time
    $\sim R/c \sim 17$\,ns (lower bound) to $R/c_s$ (upper bound)
  \item \textbf{Practical collapse time}: $\leq 1$\,ms
    (conservative upper estimate including EM dump and psi propagation)
\end{itemize}

The psi-bubble cannot be ``trapped on'' if power is removed; it is
not a persistent state. Field collapse is rapid and passive.

%=============================================================================
\section{Environmental Impact}
\label{sec:environment}
%=============================================================================

\subsection{Propellant and Exhaust}

The DFD psi-bubble creates no exhaust and consumes no propellant.
This eliminates the primary environmental hazard of conventional
rocket propulsion (combustion products, ozone layer impact from
water vapour at altitude, perchlorate contamination from solid
propellant). Environmental score: vastly superior to all existing
propulsion systems.

\subsection{Nuclear Waste and Cryogen Handling}

\begin{itemize}
  \item \textbf{Spent nuclear fuel}: After vehicle retirement, the
    reactor core constitutes spent nuclear fuel requiring geological
    repository disposal per standard nuclear waste regulations.
    Volume is small ($\sim 1$\,m$^3$ of fuel assembly) but high-level
    waste protocols apply. Design must facilitate core removal without
    crew radiation exposure.
  \item \textbf{Activated structural materials}: Materials within
    2\,m of reactor become activated. Disposal as radioactive waste.
    Low-activation alloy selection minimises this.
  \item \textbf{Tritium}: Some SMR designs produce tritium (especially
    those using $^6$Li or heavy water moderator). Assess tritium
    inventory; if $> 10^{15}$\,Bq, requires NRC licensing.
  \item \textbf{Helium cryogen}: He-4 released to atmosphere on
    venting; benign, but represents loss of rare resource. Closed-loop
    recovery system required (He recovery efficiency $> 95\%$).
\end{itemize}

\subsection{Sonic Boom Analysis}

In atmospheric operation, the psi-bubble modifies the local refractive
index for EM radiation \textit{and} acoustic propagation. The question
is whether the psi-bubble suppresses or enhances sonic boom:

\begin{itemize}
  \item \textbf{DFD psi-field and acoustic waves}: In DFD, $n_{\rm EM}
    = e^\psi$, but the speed of sound depends on different material
    properties. The psi-field is primarily an EM refractive effect;
    acoustic wave speed is set by fluid density and compressibility,
    which are not directly modified by $\psi$ (unless psi couples to
    matter density, which is the source term in DFD -- this requires
    detailed analysis of psi backreaction on matter sector).
  \item \textbf{Conventional sonic boom mechanism}: Pressure waves from
    vehicle boundary layer, projected as Mach cone.
  \item \textbf{DFD modification}: If the psi-bubble creates a smooth
    ``density gradient'' shell in the surrounding air (via EM-induced
    local heating? The mechanism here requires specification), sonic
    boom could be suppressed. However this is speculative at TRL 2;
    the conservative design assumption is: \textit{sonic boom is not
    suppressed by psi-field}. Vehicle must comply with FAA sonic boom
    regulations during atmospheric transit above Mach 1.
  \item \textbf{Practical note}: A nuclear-powered psi-bubble vehicle
    is unlikely to operate at Mach 1 in commercial airspace. Atmospheric
    transit will likely be on designated test corridors or above
    commercial airspace ($> 60,000$\,ft).
\end{itemize}

\subsection{Gravitational Wake}

The vehicle modifies the gravitational potential locally while the
psi-bubble is active. After the bubble collapses, is there a
``gravitational wake'' that affects nearby structures?

\begin{itemize}
  \item The psi-field perturbation from a 10\,m vehicle at 1\,g
    acceleration is $\delta\psi \sim a_{\rm drive} R^2 / (c^2 \cdot r^2)$
    at distance $r$
  \item At $r = 100$\,m: $\delta\psi \sim 10^{-13}$; gravitational
    acceleration perturbation $\sim 10^{-13} \cdot c^2 / R \sim 10^{-4}$\,m/s$^2$
  \item This is $10^{-5}$\,g: below the threshold for any structural effect
    on buildings or human perception
  \item No gravitational wake concern for ground operations or populated areas
\end{itemize}

%=============================================================================
\section{Regulatory Framework}
\label{sec:regulatory}
%=============================================================================

\subsection{Classification Challenge}

The DFD psi-bubble vehicle fits no existing certification category:
\begin{itemize}
  \item Not an ``aircraft'' (no lift surface, operates outside atmosphere)
  \item Not a ``launch vehicle'' (no propellant, no exhaust)
  \item Not a ``spacecraft'' in the conventional sense (nuclear-powered,
    reusable, potentially operates throughout the atmosphere)
  \item Not a ``nuclear facility'' (mobile, airborne)
\end{itemize}

A new regulatory category is required. Based on the nearest analogues,
the following framework applies:

\subsection{FAA / EASA (Atmospheric Operations)}

\begin{itemize}
  \item \textbf{Applicable framework}: FAA 14 CFR Part 460
    (Human Space Flight Requirements) for crewed operations above
    60,000\,ft; FAA 14 CFR Part 101 for uncrewed
  \item \textbf{Sonic boom}: 14 CFR Part 91 Subchapter F prohibits
    supersonic flight over populated areas; this applies during
    atmospheric transit below 60,000\,ft
  \item \textbf{Vehicle certification path}: Novel category under
    FAA Order 8110.4C; requires pre-application meeting and
    issue papers for each novel feature
  \item \textbf{Recommended approach}: Work with FAA Office of
    Commercial Space Transportation (AST) from the beginning;
    seek experimental permit (14 CFR Part 437) for development
    phase; negotiate special condition (SC) for nuclear propulsion
\end{itemize}

\subsection{Nuclear Regulatory (US NRC / DOE)}

\begin{itemize}
  \item \textbf{NRC jurisdiction}: Reactors operating within US
    territory fall under 10 CFR Part 50 (operating license) or
    10 CFR Part 52 (combined construction and operating license)
  \item \textbf{Nuclear-powered aircraft precedent}: SNAP-10A
    (space nuclear power, 1965); NTP (Nuclear Thermal Propulsion,
    NASA/DARPA ongoing) --- both provide precedent for
    non-stationary nuclear systems
  \item \textbf{Launch approval}: FAA/AST coordinates with NRC
    for nuclear-powered launch vehicles; requires environmental
    impact statement (NEPA) and nuclear risk analysis
  \item \textbf{Isotope inventory}: Core must be characterised for
    NRC regulatory purposes; SNAP-derived safety analysis required
  \item \textbf{International}: IAEA Safety Standards for nuclear
    propulsion in space (Safety Reports Series No.~40)
\end{itemize}

\subsection{Electromagnetic Compatibility (EMC)}

\begin{itemize}
  \item \textbf{Primary concern}: GJ-class SRF system at 1.3\,GHz
    is a high-power intentional radiator. In ground operation,
    must comply with FCC Part 18 (Industrial, Scientific, Medical
    equipment) or obtain experimental licence under FCC Part 5
  \item \textbf{Airborne interference}: The vehicle EM leakage (see
    Section~\ref{sec:radiation}) could interfere with GNSS receivers,
    air traffic control radar, and communication systems at ranges
    $> 1$\,km. Required: FCC waiver and coordination with FAA
    ATC spectrum management
  \item \textbf{HEMP risk}: A full-energy quench (worst-case scenario)
    with uncontrolled EM radiation into free space produces a localized
    High-altitude Electromagnetic Pulse (HEMP)-like event. This must
    be bounded and classified for military coordination
  \item \textbf{Mitigation}: Outer Faraday enclosure on vehicle
    suppresses leakage to $< -60$\,dBc of stored energy at
    1\,km distance
\end{itemize}

\subsection{Export Control and ITAR}

\begin{itemize}
  \item \textbf{ITAR classification}: Nuclear propulsion + advanced
    propulsion technology places this firmly under ITAR Category XV
    (Spacecraft Systems and Related Articles) and Category VI
    (Surface Vessels of War and Special Naval Equipment --- for
    underwater variant)
  \item \textbf{EAR dual-use}: Superconducting shell technology;
    high-power RF systems; precision cryogenics --- all have CCL
    entries under EAR. Export licence required for any technology
    transfer
  \item \textbf{Psi-field technology specifically}: Novel technology
    with no existing export control category; likely classified under
    ``600 series'' ECCN for military end-use; CJ (Commodity
    Jurisdiction) determination required from State Department
  \item \textbf{Weapons Applicability Determination}: Nuclear-powered,
    unlimited-delta-v vehicle is inherently dual-use. DOD/DOE weapons
    applicability review required before any international disclosure
  \item \textbf{Recommendation}: Engage National Security Council
    staff early; seek Fundamental Research Exclusion (FRE) for
    academic research phase; protect engineering details as
    proprietary until CJ determination complete
\end{itemize}

%=============================================================================
\section{Abort and Emergency Escape Systems}
\label{sec:abort}
%=============================================================================

\subsection{Abort Scenarios by Mission Phase}

\begin{longtable}{p{0.20\tw}p{0.25\tw}p{0.45\tw}}
\toprule
\textbf{Phase} & \textbf{Trigger} & \textbf{Primary Abort Action} \\
\midrule
\endhead
Ground, pre-launch & Any safety interlock trip & Power down; crew evacuation
  via ground egress; fire suppression activation \\
Atmospheric ascent (< 20 km) & Propulsion loss; control loss & Aerodynamic
  glide + retro-rockets; parachute if velocity $< 200$\,m/s; crew capsule
  ejection \\
High altitude (20--100 km) & Same & Retro-rocket only (air too thin for
  aerodynamic glide) + parachute deployment at $< 15$\,km \\
Orbit / deep space & Loss of propulsion & Vehicle enters elliptical drift;
  emergency life support extends to 72\,h; recovery vehicle rendezvous \\
Near massive body & Uncontrolled psi asymmetry & Immediate symmetric dump;
  retro-rocket for gravity escape; emergency psi restart at reduced power \\
\bottomrule
\end{longtable}

\subsection{Crew Escape System}

\begin{itemize}
  \item \textbf{Escape capsule}: Self-contained pressure vessel,
    independent life support 72\,h, independent retro-rocket
    ($\Delta v = 200$\,m/s), parachute + airbag landing
  \item \textbf{Ejection envelope}: Can separate from vehicle at
    any altitude from ground to orbit; works in atmosphere (aerodynamic
    drag deceleration) and space (retro-rockets)
  \item \textbf{Radiation shielding}: Capsule shielded against reactor
    accident radiation dose; crew dose $< 250$\,mGy for 72-hour stay
    after LOCA
  \item \textbf{Initiation}: Manual (crew) or automatic (triggered by
    quench protection activation, structural failure sensor, or
    loss of cabin pressure); $< 500$\,ms from initiation to separation
  \item \textbf{Vehicle separation distance}: Capsule retro-rockets
    provide $\geq 200$\,m separation within 5\,s, placing crew
    beyond structural fragment range
\end{itemize}

\subsection{Controlled Crash Landing Procedures}

If power is lost in atmosphere at low altitude (insufficient time
for crew capsule ejection):

\begin{enumerate}
  \item Immediate: All non-essential loads shed; all available power
    to retro-rockets
  \item $t + 0$--$5$\,s: Vehicle auto-orients to maximise retro-thrust
    vector against gravity
  \item $t + 5$\,s: If altitude $> 5000$\,m, crew capsule ejects
  \item $t + 5$\,s: If altitude $< 5000$\,m, \textit{vehicle} deploys
    drag surface (variable-geometry delta-wing panels folded for normal
    flight, deployed for emergency landing)
  \item $t + 30$\,s: If below $3000$\,m, vehicle parachute system deploys
    (main + drogue, rated for vehicle mass)
  \item $t + 60$\,s: Airbag impact attenuators inflate to protect reactor
    containment on ground contact
  \item On ground: Reactor emergency SCRAM confirmed; containment
    integrity check; emergency locator transmitter activates
\end{enumerate}

\subsection{Emergency Locator and Communication}

\begin{itemize}
  \item 406\,MHz ELT (Emergency Locator Transmitter), Cospas-Sarsat
    compatible, activated on impact or manual trigger
  \item Backup HF communication: 5--30\,MHz emergency frequencies,
    independent of main communication system
  \item GNSS tracking: Independent from main navigation; battery backup
    $> 120$\,h
  \item Crew survival kit: 7-day supply, designed for polar/maritime
    environments given likely test sites
\end{itemize}

%=============================================================================
\section{Consolidated Risk Matrix}
\label{sec:risk}
%=============================================================================

\begin{longtable}{p{0.06\tw}p{0.22\tw}p{0.06\tw}p{0.06\tw}p{0.06\tw}p{0.06\tw}p{0.06\tw}p{0.30\tw}}
\toprule
\textbf{ID} & \textbf{Failure Mode} &
\textbf{Sev} & \textbf{Prob} & \textbf{Det} &
\textbf{RPN} & \textbf{Risk} & \textbf{Key Mitigation} \\
\midrule
\endhead
FM-01 & Superconductor quench &
  \Scrit{10} & 6 & 8 &
  \Smed{40$^*$} & MED &
  3-layer RF+heater+dump system \\
FM-02 & Ferrite fracture &
  \Scrit{9} & 4 & 6 &
  \Smed{60$^*$} & MED &
  Composite tile + containment vessel \\
FM-03 & Cryogenic failure &
  9 & 2 & 9 &
  \Slow{36} & LOW &
  2N cryoplant + He reserve \\
FM-04 & RF feed failure &
  7 & 4 & 9 &
  \Slow{56} & LOW &
  N+1 RF chains + circulator + dump \\
FM-05 & Control failure &
  7 & 3 & 8 &
  \Slow{63} & LOW &
  TMR computing + fail-safe default \\
FM-06 & Reactor SCRAM &
  8 & 3 & 10 &
  \Slow{24} & LOW &
  Flywheel UPS + passive safe reactor \\
FM-07 & Asymmetric quench &
  \Scrit{10} & 4 & 7 &
  \Smed{20$^*$} & MED &
  Sector-resolved detection + full dump \\
FM-08 & Power bus failure &
  8 & 2 & 9 &
  \Slow{32} & LOW &
  Dual bus + independent safety UPS \\
FM-09 & Multiple simultaneous &
  \Scrit{10} & 2 & 5 &
  \Smed{120} & MED &
  FTA; physical separation; containment \\
\bottomrule
\multicolumn{8}{l}{\small $^*$ RPN shown post-mitigation. Pre-mitigation
  FM-01 = 180, FM-07 = 160 (HIGH). Mitigation is mandatory before operation.} \\
\end{longtable}

%=============================================================================
\section{Open Safety Issues Requiring Resolution Before TRL 4}
\label{sec:open}
%=============================================================================

The following safety issues are identified as blocking for progression
to prototype testing. They require experimental data or design closure
before TRL 4 (Technology Demonstration) can be claimed:

\begin{enumerate}
  \item \textbf{Quench propagation velocity in the shell geometry}: The
    DFD shell is not a solenoid or dipole magnet; $v_q$ in a large-area
    spheroidal shell has not been measured. Required: measurement of $v_q$
    in prototype shell segments. Critical for setting quench heater spacing.

  \item \textbf{Ferrite fracture toughness under combined EM + cryogenic
    loading}: No data exists for ferrite at 4\,K under combined magnetic
    and mechanical stress at the proposed operating conditions. Required:
    cryogenic mechanical testing of candidate ferrite grades.

  \item \textbf{Psi-field effect on neutron cross-sections}: A potential
    $\pm 1\%$ $k_{\rm eff}$ variation from psi-field has been identified
    as bounding assumption (Section~\ref{sec:psi_nuclear}). Required:
    dedicated experiment measuring neutron multiplication in a sub-critical
    assembly as a function of local psi-field amplitude.

  \item \textbf{Asymmetric quench dynamics timing}: The 2\,ms window
    identified for FM-07 mitigation must be validated by simulation and
    experiment. If quench propagation is faster than modelled, the
    mitigation strategy must be revised.

  \item \textbf{EM leakage from ferrite layer during full-energy operation}:
    Order-of-magnitude estimate above needs detailed EM simulation with
    as-built ferrite properties.

  \item \textbf{Psi-field gradient uniformity inside crew compartment}:
    Mapping the internal psi-field distribution at operating conditions
    is required to confirm differential crew acceleration is within
    physiological limits.
\end{enumerate}

%=============================================================================
\section{Recommended Safety Programme}
\label{sec:programme}
%=============================================================================

\subsection{Phase Gating}

\begin{longtable}{p{0.10\tw}p{0.18\tw}p{0.62\tw}}
\toprule
\textbf{TRL} & \textbf{Safety Gate} & \textbf{Required Evidence} \\
\midrule
\endhead
3 & Laboratory demo & QPS hardware demonstrated on prototype shell segment;
  quench detected and energy dumped within 100\,ms \\
4 & Component demo & Ferrite fracture tests at 4\,K completed; FM-07 timing
  validated; psi-nuclear experiment data available \\
5 & Subscale system & Full quench protection on subscale vehicle ($\sim 1$\,m,
  kJ energy); reactor SCRAM abort tested in ground rig \\
6 & System prototype & Full-scale quench protection on engineering model;
  abort procedures rehearsed; regulatory pre-application filed \\
7 & Operational demo & All FM-01 through FM-09 mitigations demonstrated
  on flight article; crew escape system tested \\
\bottomrule
\end{longtable}

\subsection{Safety Verification Plan}

For each failure mode (FM-01 to FM-09), the following verification
activities are required before operational clearance:

\begin{itemize}[nosep]
  \item \textbf{Analysis}: Physics-based models predicting failure onset,
    propagation, and damage
  \item \textbf{Simulation}: Numerical (FEA, CFD, MHD) models validated
    against test data
  \item \textbf{Test}: Deliberate fault injection tests on test articles
    (not flight hardware) at subscale and full scale
  \item \textbf{Review}: Independent review board (IRB) sign-off
    on each mitigation strategy before progression
\end{itemize}

%=============================================================================
\section{Conclusions}
%=============================================================================

The DFD psi-bubble propulsion vehicle presents a safety challenge
unprecedented in transport engineering. The combination of:
(1)~gigajoule electromagnetic energy storage in brittle superconducting
ceramics, (2)~nuclear power, (3)~superfluid cryogenic systems, and
(4)~active gravitational field modification creates a multi-dimensional
hazard landscape where standard aerospace or nuclear safety frameworks
are insufficient alone.

The following primary conclusions are reached:

\begin{enumerate}
  \item \textbf{The quench protection system is the single most critical
    safety element.} A three-layer architecture (RF dump $<10$\,ms,
    quench heaters $\leq 100$\,ms, resistive dump bank) can maintain
    peak hot-spot temperature $\leq 300$\,K for a 1\,GJ quench, making
    the vehicle survivable.

  \item \textbf{Asymmetric quench (FM-07) is the most dangerous failure
    mode unique to psi-bubble propulsion.} It must be prevented by
    sector-resolved quench detection ($< 1$\,ms response) and immediate
    full-shell dump on any asymmetry detection. This is the primary
    driver of detection system design.

  \item \textbf{The psi field itself is not a primary external hazard.}
    Tidal forces on bystanders and structures are negligible at $> 10$\,m
    range for $\leq 1$\,g acceleration. Weaponisation via psi-field
    projection is physically implausible at vehicle-scale energy budgets.

  \item \textbf{Nuclear reactor safety is manageable with proven SMR
    technology.} The psi-nuclear coupling uncertainty (potential
    $\pm 1\%$ $k_{\rm eff}$ variation) must be bounded experimentally
    before reactor operation within the psi-bubble.

  \item \textbf{No existing regulatory category applies.} A new
    framework is required; engagement with FAA AST, NRC, and ITAR
    authorities must begin at TRL 3--4, not at operational readiness.

  \item \textbf{Six open safety issues block TRL 4 progression.}
    These are tractable with focused experimental programmes; none
    requires new theoretical breakthroughs. Estimated programme
    cost to resolve: 2--3 years, \$50--150M in dedicated safety R\&D.
\end{enumerate}

The DFD psi-bubble is not inherently unsafe. Every identified
failure mode has a credible mitigation strategy grounded in
existing engineering practice (LHC quench protection, SMR safety,
aerospace abort systems). The challenge is integration: all mitigations
must work simultaneously, under mutual interference, in an extreme
multi-hazard environment. This demands a dedicated, well-resourced
safety programme from the earliest design stages, not as an afterthought.

%=============================================================================
\appendix
\section{Acronym and Symbol Reference}
%=============================================================================

\begin{longtable}{p{0.15\tw}p{0.75\tw}}
\toprule
\textbf{Term} & \textbf{Definition} \\
\midrule
\endhead
DFD & Density Field Dynamics \\
EMC & Electromagnetic Compatibility \\
ELT & Emergency Locator Transmitter \\
FMEA & Failure Mode and Effects Analysis \\
FTA & Fault Tree Analysis \\
GJ & Gigajoule ($10^9$ J) \\
HEMP & High-altitude Electromagnetic Pulse \\
ICNIRP & International Commission on Non-Ionizing Radiation Protection \\
ITAR & International Traffic in Arms Regulations \\
LHC & Large Hadron Collider \\
LOCA & Loss of Coolant Accident \\
MJ & Megajoule ($10^6$ J) \\
NDE & Non-Destructive Evaluation \\
NRC & US Nuclear Regulatory Commission \\
ODH & Oxygen Deficiency Hazard \\
QPS & Quench Protection System \\
RPN & Risk Priority Number \\
SCRAM & Safety Control Rod Axe Man (reactor emergency shutdown) \\
SIL & Safety Integrity Level \\
SMR & Small Modular Reactor \\
SRF & Superconducting Radio Frequency \\
TMR & Triple-Modular Redundant \\
TRL & Technology Readiness Level \\
UPS & Uninterruptible Power Supply \\
$J_c$ & Critical current density of superconductor \\
$T_c$ & Critical temperature of superconductor \\
$v_q$ & Quench propagation velocity \\
$\psi$ & DFD scalar density field \\
$\lambda$ & EM--psi back-reaction coupling parameter \\
\bottomrule
\end{longtable}

\end{document}
